\section{Euler integrals and summation theorems}\label{sec:integrals-sums}

\subsection{The beta integral as effective data}

For named complex parameters with positive real-part margins, define
\[
 B(x,y)=\int_0^1 t^{x-1}(1-t)^{y-1}\dd t.
\]
On the positive real axis, the powers use the real logarithm, so
\[
 \abs{t^{x-1}(1-t)^{y-1}}
 =t^{\RePart x-1}(1-t)^{\RePart y-1}.
\]
If \(\RePart x,\RePart y\ge\delta>0\), endpoint tails are bounded by
\(\eta^\delta/\delta\), up to an explicit branch-independent constant.
Thus the cutoff \(\eta\), central quadrature mesh, and evaluation precision
are primitive recursive functions of the requested accuracy.

\begin{lemma}[Beta shift identities]\label{lem:beta-shift}
Under positive real-part margins,
\begin{align}
 B(x+1,y)&=\frac{x}{x+y}B(x,y),\label{eq:beta-shift-x}\\
 B(x,y+1)&=\frac{y}{x+y}B(x,y),\label{eq:beta-shift-y}\\
 B(x+n,y)&=B(x,y)\frac{\poch x n}{\poch{x+y}n}.
 \label{eq:beta-shift-n}
\end{align}
These are effective equalities of names.
\end{lemma}

\begin{proof}
The identity
\[
 1=t+(1-t)
\]
gives \(B(x,y)=B(x+1,y)+B(x,y+1)\).  Integration by parts, performed first
on \([\eta,1-\eta]\), gives
\(xB(x,y+1)=yB(x+1,y)\).  The endpoint terms tend effectively to zero by the
positive margins.  Solving the two linear equations yields
\eqref{eq:beta-shift-x} and \eqref{eq:beta-shift-y}; induction gives
\eqref{eq:beta-shift-n}.  The finite-quadrature proof of integration by parts
uses telescoping sums and a supplied derivative modulus, as in
\cref{lem:finite-substitution}.
\end{proof}

The gamma package in the repository proves an effective version of
\[
 B(x,y)=\frac{\Gamma(x)\Gamma(y)}{\Gamma(x+y)}
\]
in the positive chamber and extends it by the gamma recurrence.  The present
arguments may take either the beta integral or that gamma quotient as the
primary name.

\subsection{Euler's hypergeometric integral}

\begin{theorem}[Euler integral]\label{thm:euler-integral}
Assume rational margins
\[
 \RePart b\ge\delta,
 \qquad
 \RePart(c-b)\ge\delta,
 \qquad \delta>0,
\]
and suppose \(1-zt\) stays apart from zero and from the chosen logarithm cut
for \(0\le t\le1\).  Then
\begin{equation}\label{eq:euler-integral}
 B(b,c-b)\F a b c z
 =\int_0^1t^{b-1}(1-t)^{c-b-1}(1-zt)^{-a}\dd t.
\end{equation}
For \(\norm z\le r<1\), the proof supplies a primitive recursive error
schedule uniform in \(z\).
\end{theorem}

\begin{proof}
On \(\norm z\le r<1\), use the effective binomial expansion
\[
 (1-zt)^{-a}=\sum_{n<N}\frac{\poch a n}{n!}(zt)^n+R_N(t,z),
\]
with a uniform geometric remainder being available from
\cref{thm:disk-convergence}.  Integrating the finite head gives
\[
 \sum_{n<N}\frac{\poch a n}{n!}z^nB(b+n,c-b).
\]
By \eqref{eq:beta-shift-n}, this is
\[
 B(b,c-b)\sum_{n<N}
 \frac{\poch a n\poch b n}{\poch c n n!}z^n.
\]
The integral of the remainder is bounded by the uniform remainder times an
effective bound for the beta weight.  Endpoint cutoffs and the central mesh
are chosen first; then \(N\) is chosen so that the integrated remainder is at
most the remaining error budget.  This proves equality of names.

For a more general \(z\) satisfying the stated branch and apartness condition,
one may either estimate the same integral directly with a local polynomial
approximation to \((1-zt)^{-a}\), or continue the already proved identity
along a certified path.  The latter is \cref{thm:path-continuation}.
\end{proof}

\begin{corollary}[The \([0,\infty)\) Euler form]\label{cor:euler-infty}
Under the corresponding margins and branch data,
\[
 \F a b c z
 =\frac1{B(b,c-b)}
 \int_0^\infty
 u^{b-1}(1+u)^{a-c}(1+(1-z)u)^{-a}\dd u.
\]
\end{corollary}

\begin{proof}
In \eqref{eq:euler-integral}, put \(t=u/(1+u)\).  The finite substitution
principle applies after cutting the improper integral at a primitive
recursive bound and estimating the algebraic tail.
\end{proof}

\subsection{Incomplete beta and elliptic integrals}

\begin{corollary}[Incomplete beta]\label{cor:incomplete-beta}
For positive real-part and branch margins,
\begin{equation}\label{eq:incomplete-beta}
 B_x(p,q)=\int_0^x t^{p-1}(1-t)^{q-1}\dd t
 =\frac{x^p}{p}\F p{1-q}{p+1}{x}.
\end{equation}
\end{corollary}

\begin{proof}
Substitute \(t=xu\).  The remaining integral is Euler's integral with
\(a=1-q\), \(b=p\), \(c=p+1\), and
\(B(p,1)=1/p\).  The branch of \(x^p\) is part of the input data.
\end{proof}

\begin{corollary}[Complete elliptic integrals]\label{cor:elliptic}
For \(|k|<1\),
\begin{align}
 K(k)&=\frac\pi2\F{1/2}{1/2}{1}{k^2},\label{eq:elliptic-K}\\
 E(k)&=\frac\pi2\F{-1/2}{1/2}{1}{k^2}.
 \label{eq:elliptic-E}
\end{align}
\end{corollary}

\begin{proof}
Euler's integral with \(b=1/2\), \(c=1\) has beta factor
\(B(1/2,1/2)=\pi\).  The substitution \(t=\sin^2\theta\) gives
\[
 \frac12t^{-1/2}(1-t)^{-1/2}\dd t=\dd\theta.
\]
For \(a=1/2\) and \(a=-1/2\), the remaining factors are respectively
\((1-k^2\sin^2\theta)^{-1/2}\) and its positive square root.  In iPRA one can
instead take the \(t\)-integrals as the definitions and verify the
trigonometric substitution by finite quadrature; no existence theorem for an
arbitrary antiderivative is needed.
\end{proof}

\subsection{Finite Chu--Vandermonde}

\begin{theorem}[Chu--Vandermonde]\label{thm:vandermonde}
For \(m\in\N\), under the finite denominator guards,
\begin{equation}\label{eq:vandermonde}
 \F{-m}{b}{c}{1}=\frac{\poch{c-b}m}{\poch c m}.
\end{equation}
This is a \StatusA{} theorem: no convergence or gamma function is involved.
\end{theorem}

\begin{proof}
Let
\[
 S_m(b,c)=\sum_{n=0}^m(-1)^n\binom mn\frac{\poch b n}{\poch c n}.
\]
We prove the formula by induction on \(m\), uniformly in \(b,c\).  Pascal's
identity gives
\begin{align*}
 S_{m+1}(b,c)
 &=S_m(b,c)-\frac bc S_m(b+1,c+1).
\end{align*}
By the induction hypothesis,
\[
 S_m(b,c)=\frac{\poch{c-b}m}{\poch c m},
 \qquad
 S_m(b+1,c+1)=\frac{\poch{c-b}m}{\poch{c+1}m}.
\]
Since
\(\poch{c+1}m=\poch c m(c+m)/c\),
\begin{align*}
 S_{m+1}(b,c)
 &=\frac{\poch{c-b}m}{\poch c m}
   \left(1-\frac b{c+m}\right)\\
 &=\frac{\poch{c-b}{m+1}}{\poch c{m+1}}.
\end{align*}
After cross-multiplication this is a quantifier-free identity between finite
products and sums.  The base \(m=0\) is immediate.
\end{proof}

\begin{corollary}[MathWorld's Kummer first formula]
For \(n\in\N\),
\[
 \F{1/2+m-k}{-n}{2m+1}{1}
 =\frac{\Gamma(2m+1)\Gamma(m+1/2+k+n)}
        {\Gamma(m+1/2+k)\Gamma(2m+1+n)},
\]
whenever the finite denominators are guarded.  The gamma notation is merely a
rewriting of a quotient of two rising factorials.
\end{corollary}

\subsection{Gauss's summation theorem}

\begin{theorem}[Effective Gauss summation in the Euler chamber]\label{thm:gauss}
Assume positive rational margins for
\(\RePart b\), \(\RePart(c-b)\), and
\(\RePart(c-a-b)\).  Then
\begin{equation}\label{eq:gauss}
 \F a b c1
 =\frac{\Gamma(c)\Gamma(c-a-b)}
        {\Gamma(c-a)\Gamma(c-b)}.
\end{equation}
The equality has a primitive recursive modulus.
\end{theorem}

\begin{proof}
Set \(z=1\) in Euler's integral.  The factor
\((1-t)^{-a}\) combines with the beta weight, giving
\[
 B(b,c-b)\F a b c1=B(b,c-a-b).
\]
Use the beta--gamma identity twice and cancel the apart gamma factors.  The
endpoint margins provide the integration modulus.  Thus the proof is not an
appeal to a limiting value at the boundary; it directly evaluates an
improper integral with an explicit tail estimate.
\end{proof}

\begin{remark}[Beyond the Euler chamber]
The classical formula only requires \(\RePart(c-a-b)>0\), apart from poles.
Starting from a parameter triple in the Euler chamber, contiguous relations
shift \(a,b,c\) by integers, and symmetry interchanges \(a,b\).  A finite
chain, together with nonzero coefficient certificates, proves many additional
instances.  Full generic extension is a \StatusC{} continuation theorem in
parameter space: the path must avoid the gamma poles and the singular
hyperplanes of the recurrence.
\end{remark}

\subsection{Kummer's theorem at \(z=-1\)}

\begin{theorem}[Kummer summation]\label{thm:kummer-sum}
Under the margins needed by the transformed Gauss sum,
\begin{equation}\label{eq:kummer-sum}
 \F a b{1+a-b}{-1}
 =\frac{\Gamma(1+a-b)\Gamma(1+a/2)}
        {\Gamma(1+a)\Gamma(1+a/2-b)}.
\end{equation}
Equivalently,
\[
 \F a b{1+a-b}{-1}
 =2^{-a}\sqrt\pi\,
  \frac{\Gamma(1+a-b)}
       {\Gamma((1+a)/2)\Gamma(1+a/2-b)}.
\]
\end{theorem}

\begin{proof}
Put \(z=-1\) in \eqref{eq:kummer-q1}.  The transformed argument is \(1\),
so Gauss's theorem gives
\[
 2^{-a}\frac{\Gamma(1+a-b)\Gamma(1/2)}
 {\Gamma((1+a)/2)\Gamma(1+a/2-b)}.
\]
The gamma duplication formula
\[
 \Gamma((1+a)/2)\Gamma(1+a/2)
 =2^{-a}\sqrt\pi\,\Gamma(1+a)
\]
turns this into \eqref{eq:kummer-sum}.  The quadratic transformation is a
finite ODE certificate; the remaining analytic content is the effective
Gauss sum and the already developed gamma identity.
\end{proof}

\subsection{Theorems at \(z=1/2\)}

\begin{theorem}[Gauss's second summation theorem]\label{thm:gauss-second}
Under the corresponding margins,
\begin{equation}\label{eq:gauss-second}
 \F a b{(1+a+b)/2}{1/2}
 =\frac{\Gamma(1/2)\Gamma((1+a+b)/2)}
        {\Gamma((1+a)/2)\Gamma((1+b)/2)}.
\end{equation}
\end{theorem}

\begin{proof}
The pullback certificate of \cref{lem:pullback-certificate} verifies
\begin{equation}\label{eq:gauss-second-quad}
 \F a b{(1+a+b)/2}{z}
 =\F{a/2}{b/2}{(1+a+b)/2}{4z(1-z)}
\end{equation}
near zero.  At \(z=1/2\), the right argument is \(1\).  Gauss's theorem
applies because
\[
 \frac{1+a+b}{2}-\frac a2-\frac b2=\frac12,
\]
and gives exactly \eqref{eq:gauss-second}.  The path from zero to one half
stays in the ordinary convergence disk, so no extra branch choice is needed.
\end{proof}

\begin{theorem}[Bailey's theorem]\label{thm:bailey}
In an initial Euler chamber, and elsewhere by certified parameter
continuation,
\begin{equation}\label{eq:bailey}
 \F a{1-a}c{1/2}
 =\frac{\Gamma(c/2)\Gamma((c+1)/2)}
        {\Gamma((c+a)/2)\Gamma((c+1-a)/2)}.
\end{equation}
\end{theorem}

\begin{proof}
Use Euler's integral with the second numerator parameter \(b=1-a\):
\[
 B(1-a,c+a-1)F
 =\int_0^1t^{-a}(1-t)^{c+a-2}(1-t/2)^{-a}\dd t.
\]
Put \(u=1-t\).  Since
\(t(1-t/2)=(1-u^2)/2\), the integral becomes
\[
 2^a\int_0^1u^{c+a-2}(1-u^2)^{-a}\dd u.
\]
Then put \(v=u^2\), obtaining
\[
 2^{a-1}B((c+a-1)/2,1-a).
\]
Divide by the original beta factor and rewrite with gamma functions:
\[
 F=2^{a-1}
 \frac{\Gamma(c)\Gamma((c+a-1)/2)}
      {\Gamma(c+a-1)\Gamma((c+1-a)/2)}.
\]
Apply the duplication formula once to \(\Gamma(c+a-1)\) and once to
\(\Gamma(c)\).  The powers of two and \(\sqrt\pi\) cancel, leaving
\eqref{eq:bailey}.  Every substitution is covered by
\cref{lem:finite-substitution}; all endpoint exponents have supplied positive
margins.
\end{proof}


\subsection{A further MathWorld value at \(z=-1\)}

\begin{theorem}[The \(\F1{-a}a{-1}\) evaluation]\label{thm:minus-one-extra}
In the initial chamber \(\RePart a>1\), and elsewhere by certified parameter
continuation,
\begin{equation}\label{eq:minus-one-extra}
 \F1{-a}a{-1}
 =1+\frac{\sqrt\pi}{2}\frac{\Gamma(a+1)}{\Gamma(a+1/2)}.
\end{equation}
\end{theorem}

\begin{proof}
Use symmetry and Euler's integral with numerator parameters \((-a,1)\):
\[
 F=(a-1)\int_0^1(1-t)^{a-2}(1+t)^a\dd t.
\]
Put \(t=2u-1\).  Then
\[
 F=(a-1)2^{2a-1}\int_{1/2}^1u^a(1-u)^{a-2}\dd u.
\]
Let \(J\) denote this last integral and let \(K\) be the corresponding
integral from \(0\) to \(1/2\).  Clearly
\(J+K=B(a+1,a-1)\).  Substituting \(u\mapsto1-u\) in \(J\) gives
\[
 J-K=\int_0^{1/2}[u(1-u)]^{a-2}(1-2u)\dd u
 =\frac{4^{1-a}}{a-1},
\]
because the integrand is
\((a-1)^{-1}\frac{d}{du}[u(1-u)]^{a-1}\).  Hence
\[
 J=\frac12\left(B(a+1,a-1)+\frac{4^{1-a}}{a-1}\right).
\]
After multiplication by \((a-1)2^{2a-1}\), the second term is \(1\), while
the first is
\[
 2^{2a-2}\frac{\Gamma(a+1)\Gamma(a)}{\Gamma(2a)}
 =\frac{\sqrt\pi}{2}\frac{\Gamma(a+1)}{\Gamma(a+1/2)}
\]
by duplication.  Every integral substitution has an effective endpoint
margin in the stated chamber.
\end{proof}

\subsection{The isolated MathWorld values}

MathWorld lists the following evaluations:
\begin{align*}
 \F{1/3}{2/3}{5/6}{27/32}&=\frac85,\\
 \F{1/4}{1/2}{3/4}{80/81}&=\frac95,\\
 \F{1/8}{3/8}{1/2}{2400/2401}&=\frac{2\sqrt7}{3},\\
 \F{1/6}{1/3}{1/2}{25/27}&=\frac{3\sqrt3}{4},\\
 \F{1/6}{1/2}{2/3}{125/128}&=\frac43\,2^{1/6},\\
 \F{1/12}{5/12}{1/2}{1323/1331}&=\frac34\,11^{1/4},
\end{align*}
and a further value at \(121/125\) involving
\(\Gamma(1/3)^3/\Gamma(1/4)^2\).

\begin{audit}[Status of the isolated values]\label{audit:isolated-values}
The archive numerically validates the displayed values to high precision and
represents every radical target by an isolated algebraic name.  A completed
iPRA proof of each algebraic value would consist of either:
\begin{enumerate}
\item a finite algebraic pullback certificate reducing the value to a
normalized point already covered above; or
\item a modular-equation certificate, an AGM/modular transformation chain,
and explicit convergence bounds.
\end{enumerate}
The source pages attribute these values to work on special hypergeometric
values, but do not print the transformation certificates.  This package does
not reconstruct all of those modular proofs, so the six displayed isolated
values and the gamma quotient are marked \StatusD{}, not silently counted as
proved.  This is the largest conspicuous omission from the finite list of
formulas on the two pages.
\end{audit}
